% Replacement for Sections II-B, II-C and II-D.
%
% Three jobs, per the T1.3 decision (Option A -- rename to match the hardware):
%   1. Compress the standard-CORDIC tutorial and attribute UD-CORDIC clearly to
%      Garrido et al. Counter-intuitively this HELPS the novelty case: the
%      submitted version reads as a tutorial on [14], which invites exactly the
%      conclusion R3 reached. Marking the background as background lets
%      Section III stand out as ours.
%   2. Make the Binary UD digit mapping unambiguous (R4-5), including the
%      convergence range, which the submitted version never stated.
%   3. Replace the "CSD" framing with what the hardware actually implements,
%      and correct Equation 8 (R4-6).
%
% TITLE CHANGE REQUIRED: the current title says "CSD-Based". Suggested
% replacement: "A Uniformly Distributed CORDIC RoPE Accelerator for LLMs".
% Every occurrence of "CSD UD-CORDIC" becomes "merged UD-CORDIC" (M-UD).

\subsection{Standard CORDIC}
\label{sec:stdcordic}

CORDIC decomposes a rotation into micro-rotations by the fixed angle set
$\arctan(2^{-i})$, each realised as a shift and an add~\cite{volder}. Because
those angles are not related by a simple recurrence, the direction $\delta_i$
of each micro-rotation must be derived from the residual angle accumulated so
far, which requires a dedicated angle datapath --- the \emph{Z-path} --- running
in lockstep with the data. The Z-path contributes adders, registers and control
in proportion to the pipeline depth, and it is the dominant overhead this work
seeks to remove.

\subsection{Binary UD-CORDIC}
\label{sec:binaryud}

Uniformly distributed (UD) CORDIC~\cite{garrido2025ud} replaces the arc-tangent
angle set with a uniform one, $\alpha_i = 2^{-i}$, so that the micro-rotation
directions can be read directly from the binary representation of the target
angle instead of being computed iteratively. The Z-path disappears entirely.
We take this construction from~\cite{garrido2025ud} unchanged; our contribution
begins in Section~\ref{sec:encoder}.

Stage $i$ of the datapath computes
\begin{equation}
  x_{i+1} = x_i - \delta_i\,(y_i \gg (i{+}1)), \qquad
  y_{i+1} = y_i + \delta_i\,(x_i \gg (i{+}1)),
  \label{eq:binstage}
\end{equation}
with $\delta_i \in \{-1, +1\}$ taken from bit $i$ of the encoded angle:

\begin{table}[h]
\centering
\caption{Binary UD digit mapping. There is no zero digit: every stage rotates.}
\label{tab:bitmap}
\begin{tabular}{ccc}
\toprule
bit $b_i$ & $\delta_i$ & effect of stage $i$ \\
\midrule
1 & $+1$ & rotate by $+\arctan(2^{-(i+1)})$ \\
0 & $-1$ & rotate by $-\arctan(2^{-(i+1)})$ \\
\bottomrule
\end{tabular}
\end{table}

Excluding $\delta_i = 0$ is not an oversight but the property that makes the
scale factor angle-independent, which Section~\ref{sec:gain} shows is what
allows RoPE to absorb it for free. The cost is that the reachable rotation is
bounded by
\begin{equation}
  |\theta| \le \sum_{i=0}^{N-1}\arctan(2^{-(i+1)}) \to 54.9^{\circ},
  \label{eq:binrange}
\end{equation}
so a target angle must first be folded into that range --- a requirement the
RoPE setting makes unavoidable, since $\theta = m\Theta_i$ grows without bound
in the token position $m$. Section~\ref{sec:frontend} handles this.

\subsection{Merged Radix-4 UD-CORDIC}
\label{sec:merged}

An $N$-bit angle needs $N$ sequential adders in the form
of~\eqref{eq:binstage}. Two adjacent binary positions can instead be handled by
a single stage operating at radix~4, halving the pipeline depth. Stage $i$ then
applies
\begin{equation}
  x_{i+1} = x_i - t_i\,y_i, \qquad y_{i+1} = y_i + t_i\,x_i,
  \qquad t_i = v_i\,2^{-2i},
  \label{eq:mergedstage}
\end{equation}
and the realised rotation is
\begin{equation}
  \theta \;\approx\; \sum_{i=0}^{N-1} \arctan\!\left(v_i\,2^{-2i}\right)
  \;\approx\; \sum_{i=0}^{N-1} v_i\,2^{-2i},
  \qquad v_i \in \{0, \pm 1, \pm \tfrac{3}{2}\},
  \label{eq:merged}
\end{equation}
the second approximation holding to better than $0.1\%$ for $i \ge 1$. The
digit $\pm\tfrac{3}{2}$ is formed as $(x \gg 2i) + (x \gg 2i{+}1)$, one adder;
$\pm 1$ is a single shift and $0$ is a bypass.

This digit set is neither canonical signed digit nor radix-4 Booth, and we
therefore do not describe the design as CSD-based. Canonical signed digit
requires digits in $\{0, \pm 1\}$ with no two non-zero digits adjacent, and
Booth recoding yields $\{0, \pm 1, \pm 2\}$; \eqref{eq:merged} is a distinct,
redundant set arising from what the merged datapath can build with one adder.
It is nonetheless an adequate number system for the purpose: the largest gap in
the digit set is $1$, so the worst residual after each stage is
$\tfrac{1}{2}\,2^{-2i}$, against a remaining capacity of
$\sum_{j>i}\tfrac{3}{2}\,2^{-2j} = \tfrac{1}{2}\,2^{-2i}$. Convergence is
therefore marginal rather than comfortable, and the reachable angle set
contains small gaps --- measured at $0.22^{\circ}$ worst case for $N = 5$ ---
but this is well inside the datapath quantisation at every word length
considered here. We verify empirically that the $N/2$-stage merged pipeline
matches the $N$-stage binary one: $0.043^{\circ}$ RMS angular error at $N = 5$
against $0.063^{\circ}$ at $N = 10$ (Table~\ref{tab:algo}).

We also evaluated the textbook $\{0, \pm 1, \pm 2\}$ set, for which $\pm 2$ is a
wire rather than an adder. It converges less marginally, but is marginally
\emph{less} accurate in practice ($0.065^{\circ}$ at $N = 5$) and, as
Section~\ref{sec:gain} shows, widens the scale-factor spread from $93\%$ to
$152\%$. We retain $\pm\tfrac{3}{2}$ on that basis.

The merged pipeline reaches $\pm 84.0^{\circ}$, wider than the binary variant,
but still short of the $360^{\circ}$ that RoPE requires.
